\section{Vector V3: Geographic Definition Gaming}
\label{sec:v3}

\subsection{Attack Description}

Before the algorithm runs, actors must make geographic definition choices:
What is the basic unit of the map?
Should water-bounded regions be connected?
How should island communities be handled?
The Census Bureau offers multiple geographic resolutions (tracts, block groups,
blocks); each produces a different adjacency graph and potentially different
algorithmic outputs.

The V3 attack exploits this upstream freedom: if the map-drawer can choose
the geographic resolution, adjacency definition, or island-connection rule
after observing partisan implications of each choice, they can steer the
output without touching any parameter or input data value.

The attack surfaces are:
\begin{enumerate}
\item \textbf{Resolution choice}: Blocks vs.\ block groups vs.\ tracts
  produce different graph topologies.
  A block-level graph has $\sim$11 million nodes nationally; a tract-level
  graph has $\sim$84,000.
  The bisection tree structure can differ substantially between resolutions
  in states where partisan geography is spatially granular.
\item \textbf{Water adjacency definition}: Whether two census units separated
  by water are treated as adjacent affects graph topology in coastal and
  lakefront states.
  Different adjacency definitions produce different bisection cuts in
  states such as Michigan, Florida, and New York.
\item \textbf{Island connection rule}: Communities on islands (e.g., the
  North Carolina Outer Banks, Florida Keys) must be connected to the
  mainland graph; the choice of connection point can affect which district
  they land in.
\end{enumerate}

\subsection{Maximum Achievable Shift}

V3 is the most consequential gaming vector for plan geometry.
F.3's threshold analysis~\citep{f3paper} establishes that resolution choice
can shift expected Democratic seat shares by up to 3.5 percentage points in
individual states (approximately 1.5 seats in a 14-seat state like North Carolina).
Nationally, the effect attenuates due to averaging: $\Delta D \leq 1.5$ seats
under the most extreme resolution switch (blocks vs.\ tracts for all 50 states).

Water adjacency and island connection rules produce smaller effects:
$\Delta D \leq 0.3$ seats nationally for water adjacency and $\leq 0.1$
for island connection rules.

The V3 effect is larger than V1 and V2 because it operates at the graph
topology level, not the parameter level.
A change in graph topology can fundamentally alter which communities are
adjacent---and thus which bisection cuts are geometrically feasible.

\subsection{Detection Mechanism}

The primary defense against V3 gaming is \emph{pre-specification} rather than
cryptographic detection.
The \dia{} should specify:
\begin{itemize}
\item \textbf{Resolution}: Census tract (GEOID suffix ``00'') as mandated
  by the P.L.\ 94-171 redistricting file.
  Block-level redistricting is not permitted without legislative approval.
\item \textbf{Water adjacency}: Tracts sharing a boundary segment of length
  $> 0$ meters in the TIGER/Line shapefile are adjacent; water-only boundary
  segments are excluded per the rule in B.1~\citep{b1paper}.
\item \textbf{Island connection}: The shortest-distance bridge rule:
  connect each isolated component to its nearest mainland node.
\end{itemize}

With these definitions fixed in statute, the adjacency hash in the manifest
provides cryptographic evidence that the correct graph was used.
Any deviation from the statutory adjacency definition changes the hash.
Detection probability: $\delta_3 = 0.95$ (5\% residual from ambiguous
boundary cases not covered by the statutory rule).
